\section{Счётчики и движение: PCNT и MCPWM}\label{les:pcnt-mcpwm}
\lessonintro
  {считать фронты аппаратно и формировать управляющий PWM}
  {уроки 4, 16, 23 и 25}
  {измеренная частота и стабильный импульс сервопривода}
  {перемычка GPIO4--GPIO5; сервопривод с отдельным питанием по желанию}

\subsection{Периферия считает вместо ISR}
При частых фронтах обработчик на каждое событие отнимает CPU и может пропустить импульс.
PCNT считает фронты аппаратно; процессор периодически читает результат. Канал умеет
учитывать направление по второй линии, поэтому квадратурный энкодер можно декодировать аппаратно.
Фильтр коротких импульсов помогает с помехами, но не заменяет устранение дребезга механической
кнопки: полезный сигнал нельзя удалить вместе с шумом.

\subsection{Частотомер на одной плате}
Отключите LED и кнопку от GPIO4/5 и соедините GPIO перемычкой. LEDC формирует 1000 Гц,
PCNT считает положительные фронты. За секунду ожидаем около 1000 событий; точное время окна
измеряем через \code{millis()}. Все логические сигналы в опыте имеют уровень 3,3 В.

\begin{figure}[H]
\centering
\begin{circuitikz}[scale=0.95,transform shape]
  \node[blk,fill=blkcpu,minimum width=2.8cm,minimum height=3.2cm] (m) at (0,0) {};
  \node[font=\sffamily\small\bfseries] at (0,2) {ESP32-S3};
  \node[font=\sffamily\scriptsize,anchor=east] at (1.4,0.7) {\pin{4}};
  \node[font=\sffamily\scriptsize,anchor=east] at (1.4,-0.7) {\pin{5}};
  \draw[thick,accent] (1.4,0.7) -- (3.4,0.7) -- (3.4,-0.7) -- (1.4,-0.7);
  \node[font=\sffamily\scriptsize,text=accent,anchor=west,align=left] at (3.7,0.7) {выход LEDC, \SI{1}{kHz}};
  \node[font=\sffamily\scriptsize,text=accent,anchor=west,align=left] at (3.7,-0.7) {вход PCNT};
  \node[font=\sffamily\scriptsize,text=accent,anchor=west] at (3.7,0) {перемычка (один провод)};
\end{circuitikz}
\caption{Частотомер на одной плате: перемычка соединяет выход генератора \pin{4} со входом счётчика
\pin{5}. Другие детали к этим выводам не подключайте}
\end{figure}

\codecap{PCNT: аппаратный подсчёт фронтов генератора}
\codeinput{l31_pcnt}

Граница счётчика конечна. При достижении верхнего предела 30000 аппаратный счётчик
обнуляется. В примере разность корректируется по модулю 30000: за окно должно прийти меньше
30000 событий, а счёт идёт только вверх. Для реального двунаправленного энкодера нужны
watch points и расширение счётчика; этот частный приём туда напрямую не переносится.
Чтение и очистка отдельными командами могли бы потерять фронт, поэтому каждую секунду не очищаем.

\subsection{MCPWM и форма управляющего сигнала}
LEDC подходит для яркости и простого PWM. MCPWM добавляет связанные генераторы,
компараторы, синхронизацию, захват и обработку аварий. Dead time --- пауза между выключением
одного плеча моста и включением другого. Она снижает риск одновременного открытия плеч.
GPIO всё равно управляет входом драйвера, а не обмоткой или силовым ключом напрямую.

\codecap{MCPWM: период 20 мс, высокий уровень 1500 мкс}
\codeinput{l31_mcpwm_servo}

Таймер задаёт шкалу 1 мкс и период 20000 отсчётов. Событие начала периода включает выход,
совпадение с компаратором 1500 выключает его. Получается команда среднего положения
на GPIO6. Питание сервопривода отдельное, GND объединён. Допустимый диапазон импульсов
проверьте для модели: начинайте с середины, не с механических упоров.

\begin{figure}[H]
\centering
\begin{circuitikz}[scale=0.95,transform shape]
  \node[blk,fill=blkcpu,minimum width=2.8cm,minimum height=3.4cm] (m) at (0,0) {};
  \node[blk,fill=blkper,minimum width=2.8cm,minimum height=2.4cm] (d) at (8,0.3) {};
  \node[blk,fill=blkrtc,minimum width=2.8cm,minimum height=1.2cm,align=center] (s) at (8,-2.6) {};
  \node[font=\sffamily\scriptsize,align=center] at (8.5,-2.6) {Внешний\\[-2pt] источник \SI{5}{V}};
  \node[font=\sffamily\small\bfseries] at (0,2.1) {ESP32-S3};
  \node[font=\sffamily\small\bfseries,align=center] at (8,1.95) {Сервопривод};
  \node[font=\sffamily\scriptsize,anchor=east] at (1.4,0.9) {\pin{6}};
  \node[font=\sffamily\scriptsize,anchor=east] at (1.4,-1) {GND};
  \foreach \y/\n in {0.9/{сигнал},0.3/{+5 В},-0.3/GND}
    \node[font=\sffamily\scriptsize,anchor=west] at (6.6,\y) {\n};
  \foreach \y/\n in {-2.3/{+},-2.9/{$-$}}
    \node[font=\sffamily\scriptsize,anchor=west] at (6.6,\y) {\n};
  \draw[thick,warn]     (1.4,0.9) -- (6.6,0.9);
  \draw[thick,esp]      (6.6,0.3) -- (5.6,0.3) -- (5.6,-2.3) -- (6.6,-2.3);
  \draw[thick,black!70] (6.6,-0.3) -- (4.8,-0.3) -- (4.8,-2.9) -- (6.6,-2.9);
  \draw[thick,black!70] (1.4,-1) -- (4.8,-1);
  \fill[black!70] (4.8,-1) circle (1.5pt);
  \node[font=\sffamily\scriptsize,text=black!70,anchor=north] at (3.1,-1) {общий GND};
\end{circuitikz}
\caption{Подключение сервопривода: сигнал с \pin{6}, питание \SI{5}{V} от отдельного источника, земли платы
и источника соединены. Провод питания сервопривода к 3V3 платы не подключайте}
\end{figure}

\subsection{Первый контур обратной связи}
Управление скоростью состоит из задания, измерения PCNT, ошибки и изменения PWM.
Проверьте каждый блок отдельно. Перевод числа событий в обороты зависит от импульсов на оборот
и выбранных фронтов: квадратурный режим x4 даёт иной коэффициент, чем одна линия и один фронт.
Период регулирования должен быть известен; случайные задержки ухудшают поведение контура.

\begin{trouble}
Ноль событий: нет перемычки или выбран другой GPIO. Удвоенная частота: считаются оба фронта.
Потеря положения: неверный коэффициент энкодера или не обработана граница счётчика.
\end{trouble}
\begin{summary}
PCNT считает события, MCPWM задаёт связанные выходы. Для измерения нужны время окна,
обработка границ и правильное число фронтов на оборот.
\end{summary}

\begin{selfcheck}
\begin{enumerate}
\item Сколько положительных фронтов даёт 1000 Гц за 250 мс?
\item Зачем обрабатывать границы PCNT?
\item Что означает dead time?
\end{enumerate}
\end{selfcheck}
\needspace{20\baselineskip}
\begin{tasks}
\begin{enumerate}
\item \lv{1} Нарисуйте соединение генератора и PCNT.
\item \lv{1} Запишите пять измерений около 1000 Гц.
\item \lv{2} Сравните 100, 1000 и 5000 Гц при двух длинах окна.
\item \lv{2} Измените компаратор на 1400 мкс и сначала проверьте анализатором.
\item \lv{3} Подключите энкодер с аппаратным определением направления и проверьте возврат к нулю.
\item \lv{4} Спроектируйте регулирование скорости низковольтного мотора с ограничением PWM и аварийной остановкой.
\end{enumerate}
\end{tasks}

\paragraph{Документация для продолжения}\href{https://docs.espressif.com/projects/esp-idf/en/v5.4/esp32s3/api-reference/peripherals/pcnt.html}{Официальное руководство по теме}.
\href{https://docs.espressif.com/projects/esp-idf/en/v5.4/esp32s3/api-reference/peripherals/mcpwm.html}{Архитектура MCPWM}.
